\chapter{CẤU HÌNH WOOCOMMERCE VÀ CÁC TÍNH NĂNG BÁN HÀNG}

Sau khi hoàn tất cài đặt WordPress và tích hợp giao diện \textbf{Sen Việt Tea}, nhóm thực hiện dự án tiến hành cấu hình hệ thống WooCommerce. WooCommerce đóng vai trò là nền tảng thương mại điện tử cốt lõi, quản lý toàn bộ chu trình mua sắm từ việc duyệt xem sản phẩm, thêm vào giỏ hàng, tính toán chi phí vận chuyển cho đến bước thanh toán hoàn tất. Chương này trình bày chi tiết các thiết lập hệ thống và tính năng bán hàng dựa trên cơ sở dữ liệu thực tế của dự án.

\section{Danh mục và Sản phẩm}

\subsection{Phân chia danh mục sản phẩm}
Việc phân chia danh mục sản phẩm một cách khoa học giúp tối ưu hóa trải nghiệm người dùng, giảm thiểu thời gian tìm kiếm và định hướng hành vi mua sắm hiệu quả. Nhóm thực hiện dự án đã phân loại danh mục sản phẩm của Sen Việt Tea thành bốn nhóm chính, phù hợp với thói quen tiêu dùng trà của người Việt:

\begin{itemize}
    \item \textbf{Trà Thái Nguyên:} Đáp ứng nhu cầu thưởng thức các dòng trà xanh thuần mộc truyền thống với vị chát dịu và hậu ngọt sâu đặc trưng.
    \item \textbf{Trà Ướp Hoa:} Bao gồm các sản phẩm trà sen Tây Hồ, trà nhài mang hương thơm thanh tao tự nhiên, thích hợp làm quà tặng hoặc dành cho người thích phong vị trà nhẹ nhàng.
    \item \textbf{Trà Đặc Sản Cổ Thụ:} Cung cấp các dòng trà cổ thụ Shan Tuyết hoang dã từ vùng núi cao Tây Bắc và các loại trà đặc sản vùng miền độc đáo, mang đến những trải nghiệm thưởng thức mới lạ cho người yêu trà.
    \item \textbf{Phụ Kiện Trà Cụ:} Cung cấp các sản phẩm phụ trợ thiết thực như ấm chén đất nung Bát Tràng, tống trà, lọc trà. Danh mục này không chỉ đáp ứng trọn vẹn nhu cầu thưởng trà chuyên nghiệp mà còn hỗ trợ chiến lược bán kèm nhằm gia tăng giá trị trung bình trên mỗi đơn hàng.
\end{itemize}

Để tối ưu hóa hiệu quả tìm kiếm tự nhiên, các đường dẫn tĩnh của từng danh mục được cấu hình ngắn gọn và tường minh, ví dụ: \texttt{/danh-muc/tra-thai-nguyen}. Điều này hỗ trợ các công cụ tìm kiếm lập chỉ mục cấu trúc trang web một cách chính xác nhất.

\subsection{Đưa sản phẩm lên hệ thống}
Nhằm tối ưu hóa thời gian và đảm bảo tính nhất quán của dữ liệu, dự án sử dụng phần mở rộng \textbf{Product Import Export for WooCommerce} để tải dữ liệu sản phẩm hàng loạt từ tệp cấu trúc CSV lên hệ thống. Công cụ này hỗ trợ tự động đồng bộ hóa các thông tin cốt lõi bao gồm hình ảnh sản phẩm, mô tả chi tiết, giá bán và mã định danh sản phẩm SKU vào cơ sở dữ liệu.

Hiện tại, hệ thống đã chuẩn hóa và đưa vào vận hành chính thức 30 sản phẩm đặc trưng của Sen Việt Tea, tiêu biểu như:
\begin{itemize}
    \item \textit{Bộ 4 Ly Trà Thủy Tinh Hai Lớp Cách Nhiệt}
    \item \textit{Ấm Trà Đất Nung Dáng Tây Thi Cổ Điển}
    \item \textit{Hũ Đựng Trà Gốm Tử Sa Nắp Kín}
    \item \textit{Lọc Trà Sứ Viền Inox Siêu Mịn}
    \item \textit{Tống Trà Thủy Tinh Chịu Nhiệt Cao Cấp}
\end{itemize}

Các mặt hàng này được thiết lập thuộc loại \textbf{Sản phẩm đơn giản} do mỗi sản phẩm có một mức giá cố định và không phân chia thành các thuộc tính biến thể như màu sắc hoặc kích thước. Đồng thời, tính năng quản lý kho hàng tự động được kích hoạt nhằm kiểm soát số lượng vật lý trong kho. Khi số lượng chạm mốc 0, hệ thống sẽ tự động vô hiệu hóa nút mua hàng để tránh tình trạng khách hàng đặt mua các sản phẩm đã hết hàng.

\section{Vận chuyển và Thuế}

\subsection{Khu vực vận chuyển và Phí giao hàng}
Theo các nghiên cứu về hành vi mua sắm trực tuyến, chi phí vận chuyển phát sinh đột ngột ở bước thanh toán là một trong những nguyên nhân hàng đầu dẫn đến việc khách hàng từ bỏ giỏ hàng. Để giải quyết vấn đề này và đơn giản hóa quy trình mua sắm, dự án đã xây dựng chính sách vận chuyển đồng giá áp dụng cho mọi khu vực địa lý trong phạm vi toàn quốc thông qua thiết lập phân vùng vận chuyển trong WooCommerce.

Bất kể địa chỉ giao nhận của khách hàng ở khu vực nào, hệ thống đều hiển thị rõ ràng hai phương thức vận chuyển lựa chọn:
\begin{enumerate}
    \item \textbf{Giao hàng tiêu chuẩn:} Có mức phí cố định là 35.000 VNĐ, thời gian vận chuyển dự kiến từ 2 đến 4 ngày làm việc.
    \item \textbf{Giao hàng nhanh hỏa tốc:} Có mức phí cố định là 200.000 VNĐ, đáp ứng nhu cầu giao nhận khẩn cấp trong ngày đối với các khu vực nội thành, phù hợp với các đơn hàng quà tặng cấp bách.
\end{enumerate}

Chính sách đồng giá giúp doanh nghiệp dễ dàng kiểm soát chi phí vận hành với các đơn vị chuyển phát đối tác, đồng thời mang lại sự minh bạch cho khách hàng ngay từ đầu chu trình thanh toán.

\subsection{Cấu hình Thuế}
Tại thị trường Việt Nam, người tiêu dùng trực tuyến thường ưu tiên việc hiển thị mức giá cuối cùng của sản phẩm thay vì phải chịu thêm các khoản phụ phí thuế phát sinh ở bước thanh toán. Do đó, hệ thống được cấu hình hiển thị \textbf{Giá đã bao gồm thuế}. Mặc dù giá niêm yết trên website là giá trọn gói, hệ thống WooCommerce vẫn thực hiện tách xuất phần thuế giá trị gia tăng trên hóa đơn điện tử nhằm đảm bảo tính chính xác về mặt kế toán và tuân thủ các quy định pháp lý hiện hành.

\section{Cổng thanh toán}

\subsection{Thanh toán khi nhận hàng}
Bất chấp sự phát triển mạnh mẽ của các phương thức thanh toán số, hình thức giao hàng thu tiền tận nơi vẫn chiếm một tỷ trọng lớn trong thương mại điện tử tại Việt Nam. Đặc biệt đối với một thương hiệu mới như Sen Việt Tea, phương thức thanh toán này giúp xây dựng niềm tin ban đầu với người tiêu dùng mới, cho phép họ nhận hàng, kiểm tra quy cách đóng gói và chất lượng sản phẩm trước khi hoàn tất giao dịch tài chính.

\subsection{Tích hợp cổng thanh toán quốc tế PayPal}
Nhằm đa dạng hóa các phương thức thanh toán số và hướng tới tệp khách hàng quốc tế, dự án tích hợp cổng thanh toán quốc tế PayPal thông qua ứng dụng mở rộng \textbf{WooCommerce PayPal Payments}. Sau khi kết nối thông tin xác thực API bao gồm mã định danh Client ID và khóa bí mật Secret Key từ tài khoản doanh nghiệp PayPal, website có khả năng chấp nhận thanh toán bảo mật từ các loại thẻ thanh toán quốc tế như Visa hoặc Mastercard, cũng như số dư tài khoản ví PayPal. Tính năng này tạo điều kiện thuận lợi cho đối tượng kiều bào nước ngoài mua quà tặng người thân tại Việt Nam hoặc khách du lịch quốc tế có nhu cầu mua trà đặc sản. Hệ thống sẽ tự động xử lý việc quy đổi tỷ giá ngoại tệ và ghi nhận chi tiết dòng tiền trong báo cáo quản trị WooCommerce.

\section{Mã giảm giá và Chương trình ưu đãi}

Để thúc đẩy quyết định mua sắm và gia tăng tỷ lệ chuyển đổi đơn hàng, dự án tận dụng công cụ quản lý mã ưu đãi tích hợp sẵn trong nhân WooCommerce thay vì cài đặt thêm các tiện ích bổ sung của bên thứ ba, qua đó tối ưu hóa hiệu năng tải trang của hệ thống. Hiện tại, trang web triển khai ba mã giảm giá chiến lược phục vụ các mục tiêu tiếp thị khác nhau:
\begin{itemize}
    \item \textbf{FREESHIP:} Miễn phí vận chuyển tiêu chuẩn trị giá 35.000 VNĐ cho các đơn hàng đủ điều kiện, thường dùng để hỗ trợ các chiến dịch quảng bá quy mô toàn quốc.
    \item \textbf{GIAM10:} Giảm 10\% trên tổng giá trị đơn hàng, áp dụng tri ân nhóm khách hàng thân thiết hoặc khách hàng mua sắm với số lượng lớn.
    \item \textbf{CHAO2026:} Khấu trừ trực tiếp 50.000 VNĐ vào tổng hóa đơn. Đây là mã ưu đãi đặc biệt chào mừng năm mới, tạo động lực tâm lý tích cực kích thích khách hàng hoàn tất giỏ hàng.
\end{itemize}

Khách hàng có thể dễ dàng áp dụng các mã ưu đãi này ngay tại giao diện giỏ hàng hoặc trang thanh toán. Hệ thống sẽ tự động đối chiếu điều kiện áp dụng, kiểm tra tính hợp lệ của mã và cập nhật lại tổng số tiền thanh toán của đơn hàng một cách nhanh chóng và chính xác.

\section*{Tiểu kết chương 4}
Chương 4 đã mô tả chi tiết toàn bộ quy trình thiết lập các luồng nghiệp vụ kinh doanh cốt lõi của website Sen Việt Tea, bao gồm danh mục trà, thuộc tính biến thể, và quy trình thanh toán khép kín. Các cơ sở dữ liệu đã được nạp thành công, tạo nên một cửa hàng trực tuyến với dữ liệu hoàn chỉnh, sẵn sàng cho công tác kiểm thử và triển khai các chiến lược tiếp thị được trình bày ở Chương 5.